% ============================================================================
\clearpage
\section*{v4.9.0-r08 Source-Only: Frame Drift, Atomic Clocks and the Hubble Tension}
\addcontentsline{toc}{section}{v4.9.0-r08 Source-Only: Frame Drift, Atomic Clocks and the Hubble Tension}
\label{sec:v490-r08-drift-clocks-hubble}
% ============================================================================

\paragraph{Purpose.}
Proposition~\ref{prop:H0-tension} explains the Hubble tension by a drift of the frame offset.  r17 and r18 tested that
drift against the expansion history.  This layer tests it against the constancy of the fine-structure constant, using
the source's own rule for which quantities stay fixed.  It also checks the other explanations of the tension available
in the source, settles the sign of \(k\) for \(\alpha\) left open in r03, names a frame-invariant clock, and states what
the drift means for H15 of r07.  Files: \texttt{llm/Erkundungsprotokoll\_Weinberg\_Kegel.md} (entries E10 to E13) and
scripts \texttt{llm/k47\_drift\_weinberg.py}, \texttt{llm/k48\_feste\_groessen.py}.  No new prediction is made.  Earlier
layers are not rewritten.

\noindent\textbf{Source-only status.}
This v4.9.0-r08 layer is source-only.  It does not build a PDF, does not assign a new DOI,
and does not authorize a public release.  The v4.9.0 Version DOI
\href{https://doi.org/10.5281/zenodo.23534954}{10.5281/zenodo.23534954} and the Concept DOI
\href{https://doi.org/10.5281/zenodo.19210728}{10.5281/zenodo.19210728} are unchanged by it.

\subsection*{v4.9.0-r08 what stays fixed under drift}
\label{subsec:v490-r08-fixed}

The source fixes this in four places.
\begin{itemize}[leftmargin=2em]
\item The definition in \S\ref{sec:phase-constants}: \(X^{(\mathrm{photon})}=\gL^{-k(X)}X^{(\mathrm{matter})}\).
\item Theorem~\ref{thm:k-values}: \(k(c)=0\), \(k(\gBI)=0\) (``frame-invariant by construction'').
\item The corollary ``\(\alpha\)-drift as a frame effect'' in the same section: ``the hardware value
\(\alpha^{(\mathrm{photon})}\) is strictly constant'', with \(\alpha^{(\mathrm{matter})}(t)=\gL^{(\mathrm{local})}(t)\,
\alpha^{(\mathrm{photon})}\).
\item Item (1) of Proposition~\ref{prop:saturation}, which writes the neutrino mass in the same way.
\end{itemize}
Photon-frame values and quantities with \(k=0\) stay fixed.  Matter-frame values drift with a local factor
\(\gL^{(\mathrm{local})}(t)^{k}\).  The structural factor \(\gL=\gBI/\gstring\) is fixed, since both inputs are fixed
numbers.  The source thus separates a fixed structural factor from a local one that may drift.  This is the separation
r18 offered.

\subsection*{v4.9.0-r08 the sign of \texorpdfstring{$k$}{k} for \texorpdfstring{$\alpha$}{alpha}}
\label{subsec:v490-r08-alpha-sign}

r03 (F2) found that Theorem~\ref{thm:k-values}, whose row reads ``\(\alpha^{-1}\ldots k=1\)'', and the corollary above are
inverse to each other, and left the choice open.  The main body uses \(k\) for \(\alpha\) in three places:
\begin{itemize}[leftmargin=2em]
\item prediction P5 among the predictions around Corollary~\ref{cor:falsifiable}, with
\(\gL^{k(\alpha)-k(G_N)}=\gL^{-1}\);
\item Proposition~\ref{prop:muon-g2}, ``\(k(\alpha_s)=1\) (like the electromagnetic coupling \(\alpha\))'' with
\(\alpha_s^{(\mathrm{photon})}=\gL^{-1}\alpha_s^{(\mathrm{matter})}\);
\item the corollary, \(\alpha^{(\mathrm{matter})}=\gL\,\alpha^{(\mathrm{photon})}\).
\end{itemize}
All three mean \(k(\alpha)=+1\) for \(\alpha\) itself.  No photon-frame value of \(\alpha^{-1}\) is used anywhere in the main
body.

\paragraph{Resolution (reading decision).}
The table row shows the CODATA value of \(\alpha^{-1}\), and its entry \(k=1\) belongs to \(\alpha\).  With this reading all
three uses agree and nothing else changes.  The resolution changes no number of the main body.

\subsection*{v4.9.0-r08 the drift against atomic clocks}
\label{subsec:v490-r08-clocks}

\paragraph{Theoremlet v4.9.0-r08.1.}
Suppose the drift of Proposition~\ref{prop:H0-tension}, \(1.6\times10^{-12}\) per year in \(\bsym\), acts through the local
frame factor, as r18 reads it.  Then the corollary gives \(|{\rm d}\ln\alpha/{\rm d}t|\) for each choice of the factor:

\begin{center}
\small
\begin{tabular}{@{}lrrr@{}}
\toprule
factor & exponent needed for \(1.083\) & \(|{\rm d}\ln\alpha/{\rm d}t|\) per year & \(|\Delta\alpha/\alpha|\) since \(z=1100\)\\
\midrule
\(\gL=\cosh2\eta\) & \(1.52\) & \(3.9\times10^{-12}\) & \(5.4\%\)\\
neighbour, \(\cosh\eta\) & \(4.72\) & \(1.3\times10^{-12}\) & \(1.7\%\)\\
neighbour, \(e^{\eta}\) & \(2.54\) & \(2.3\times10^{-12}\) & \(3.2\%\)\\
\bottomrule
\end{tabular}
\end{center}

The comparison of two \({}^{171}\mathrm{Yb}^+\) clocks gives \((1/\alpha)\,{\rm d}\alpha/{\rm d}t=(1.0\pm1.1)\times10^{-18}\)
per year (Lange et al., Phys.\ Rev.\ Lett.\ 126, 011102, 2021).  Every choice lies about six orders of magnitude above
that bound.

Independently of the reading: if a drifting factor with exponent \(k(H_0)\) produces the tension \(\ln1.083\), then
\(\alpha\) changes by \(\ln1.083\cdot k(\alpha)/k(H_0)\).  For a drift linear in time, the clocks require
\(k(\alpha)/k(H_0)<5.5\times10^{-7}\) at \(2\sigma\).

\paragraph{Proof sketch.}
\({\rm d}\ln\gL/{\rm d}t=4\bsym(1-\bsym^2)^{-2}\dot\bsym/\gL\), and correspondingly for the neighbour factors.  The
exponents follow from the ratios of the factors at \(z=0\) and \(z=1100\).  Checked numerically.  \(\square\)

\paragraph{Why measuring through photons does not help.}
A clock comparison measures the ratio of two transition frequencies of atoms.  A common frame factor on both frequencies
cancels.  The dependence on \(\alpha\) comes from inside the atom, which is matter.  If the atom's internal physics were
photon-frame physics, there would be no drifting \(\alpha^{(\mathrm{matter})}\) at all.  The corollary's annual-drift
prediction would then lapse too, and Proposition~\ref{prop:H0-tension} would lose its mechanism.

\paragraph{Consequence.}
In its frame-factor form, Proposition~\ref{prop:H0-tension} cannot explain the tension and keep \(\alpha\) constant at the
same time.  Two ways out would remain, and the source contains neither.  One is \(k(\alpha)\approx0\), which contradicts the
resolution above and the corollary.  The other is a separate factor for the expansion that does not touch the
constants.  r03 already bounded the annual modulation of the local factor by \(10^{-17}\).  The present layer applies the
same clocks to the secular drift.

\subsection*{v4.9.0-r08 other explanations in the source}
\label{subsec:v490-r08-other}

\begin{enumerate}[leftmargin=2em]
\item \emph{Static frame typing.}  The value inferred from the microwave background would be read in another frame than
the local value.  A constant factor disturbs no clock.  By Corollary~\ref{cor:gue-frame} a clock frequency transforms with
\(\gL\), so \(k(H_0)=1\), which gives a factor \(1.864\) (neighbour factors \(1.197\), \(1.854\)) against the observed
\(1.083\pm0.017\).  Only fractional exponents near \(0.13\) or \(0.44\) would fit, and Theorem~\ref{thm:k-values} gives no
basis for them.  The observed ratio is too loose to single out an exponent.
\item \emph{A change confined to the early universe.}  A larger electron mass at recombination, by about one to two
percent, is discussed in the literature as a relief of the tension.  With the drift of
Proposition~\ref{prop:H0-tension}, \(\bsym\) was smaller earlier, so \(m_e^{(\mathrm{matter})}\) with \(k=1\) was smaller.
That is the wrong direction, and the same holds for \(\alpha\) with \(k(\alpha)=+1\).  The opposite sign of the drift, the one
r17 found Theorem~\ref{thm:dark-energy} to require, would point the right way.  The change would then have to stop before
\(z\approx3\), because of quasar spectra and clocks, and at \(5\%\) it would be too large.
\end{enumerate}
No explanation of the tension within the source survives the clocks, the quasar spectra and the size of the effect
together.  Proposition~\ref{prop:H0-tension} stays a hypothesis without a working mechanism.

\subsection*{v4.9.0-r08 a frame-invariant clock}
\label{subsec:v490-r08-compton}

\paragraph{Theoremlet v4.9.0-r08.2.}
By Theorem~\ref{thm:k-values} the Compton frequency \(m_ec^2/\hbar\) has \(k=k(m_ec^2)-k(\hbar)=1-1=0\).  A clock ticking
at a Compton frequency is frame-invariant in the source's sense.  An atomic transition frequency scales as
\(\alpha^2m_ec^2/\hbar\) and has \(k=2k(\alpha)=2\).  The ratio of the two clocks reads \(\gL^{(\mathrm{local})}(t)^{\pm2}\)
directly, a local test of the frame factor that needs no cosmology.

\paragraph{Proof sketch.}
Additivity of \(k\) (Lemma v4.8.0-r03.1).  \(\square\)

\paragraph{Remark.}
Clocks tied to the Compton frequency of a mass have been built by atom interferometry (Lan et al., Science 339, 554, 2013).
A comparison of such a clock with an atomic clock would read the drift of r08.1 at twice the rate.  Existing clock
comparisons already exclude that drift.  The clock concerns local readout only.  It touches neither the local selector
guard nor the forbidden global photon-clock statement.

\subsection*{v4.9.0-r08 consequence for H15}
\label{subsec:v490-r08-h15}

H15 (r07) is stated with the structural factor \(\gL=\gBI/\gstring\), which is fixed.  It predicts a weak mixing angle that
is constant in time.  In addition, \(m_W/m_Z\) is a ratio of masses, and if both carry the same \(k\) the frame factor
cancels.  A drift of \(\theta_W\), with \(\sin^2\theta_W=0.242\) at \(z=1100\), would follow only if H15 were read with the
local factor, which the source's rules do not suggest.  The pre-registered comparison of H15 is unaffected.

\subsection*{v4.9.0-r08 register entries}
\label{subsec:v490-r08-register}

\begin{center}
\small
\begin{tabular}{@{}>{\raggedright\arraybackslash}p{0.34\textwidth}>{\raggedright\arraybackslash}p{0.58\textwidth}@{}}
\toprule
statement & status after v4.9.0-r08\\
\midrule
Prop.~\ref{prop:H0-tension} & frame-factor drift incompatible with atomic clocks by about six orders of magnitude in every reading (r08.1); static typing and early-only change do not work; hypothesis without working mechanism\\
Thm.~\ref{thm:k-values}, row \(\alpha^{-1}\) & resolved: \(k=1\) belongs to \(\alpha\); consistent with P5, Prop.~\ref{prop:muon-g2} and the corollary (closes r03 F2)\\
fixed quantities under drift & photon-frame values and \(k=0\); structural \(\gL=\gBI/\gstring\) fixed; local factor may drift\\
Compton clock & \(k=0\), frame-invariant; ratio to an atomic clock reads \(\gL^{(\mathrm{local})}{}^{\pm2}\) (r08.2)\\
H15 & predicts a constant \(\theta_W\); comparison unaffected\\
\bottomrule
\end{tabular}
\end{center}

\subsection*{v4.9.0-r08 ledger and guard update}
\label{subsec:v490-r08-ledger-guard}

\[
R_{\rm drift}^{v490\text{-}r08}=(0,0,1)
\]
for (drift against clocks decided, sign of \(k(\alpha)\) resolved, a mechanism for the Hubble tension open).

The layer spends no guard:
\[
R_{\rm glob}^{(v490\text{-}r08,+)}=(0,1,1,1,1),
\qquad
\Gamma_{\rm guard}^{v490\text{-}r08}=(1,1,1,1,1),
\qquad
\beta_4=1/256\ \text{unspent}.
\]
It does not imply exact finite \(C5\) zero closure, exact all-readout terminality,
\(C6/T2\), a global photon-clock/RH theorem, or global all-readout closure.

